%==================================================================
\section{The 2022 Mid-Semester Paper}

\setcounter{pq}{0}

\begin{center}
\fbox{\parbox{0.93\textwidth}{\centering
\textbf{INDIAN STATISTICAL INSTITUTE}\\
\textbf{Mid Semestral Examination}\\
B.Stat.\ --- 3rd Year (Semester V)\\
\textit{Design and Analysis of Algorithms}\\[3pt]
Date: September 22, 2022 \qquad Maximum Marks: 60 \qquad Duration: 2:30 Hours\\[3pt]
\small\textit{Note: You may answer any part of any question,\\ but maximum you can score is 60.}
}}
\end{center}

\qhead{5+5}
\pt{a} $A$ is an unsorted array of size $n$ where $A[i]\in\{1,2,3,4\}$ for each
$i\in\{0,1,2,\dots,n-1\}$. Write an algorithm to sort the array $A$ in $\Oh n$ time.

\pt{b} Let $S=\langle(a_i,b_i):a\in\N,\ b\in\{0,1\}\rangle$ be a sequence stored in an array of
size $2n$. Write a linear-time algorithm using $\Oh1$ space that rearranges $S$ such that all
the tuples $(a_i,b_i)$ having $b_i=0$ appear at the beginning.

\qhead{7+3}
\pt{a} Let $S$ be a point set in the plane where $|S|=n$. Design an algorithm to find the convex
hull of $S$ in $\Oh{n\log n}$ time.

\pt{b} Prove that the lower bound of an algorithm that finds the convex hull of a given set of
points $S$ is $\Om{n\log n}$, where $|S|=n$.

\qhead{5+5}
\pt{a} An inversion in an array $A$ is a pair $(i,j)$ having the property $i<j$ and $A[i]>A[j]$.
Let $X$ be the total number of inversions in an array. Calculate $\E(X)$.

\pt{b} Let $T(n)$ be a positive and non-decreasing function such that $T(n)=a\,T(n/b)+c\,n^k$
and $T(1)=d$ where $a,c>0$, $k,d\ge0$, $b\ge2$. Prove that
$T(n)=\Th{n^{\log_b a}}$.\footnote{The scanned paper carries a handwritten addition
``\emph{and} $a>b^k$'' --- without it the claim is false (see \S2), so read the condition as
part of the question.}

\qhead{10}
Let $S=\langle a_1,a_2,\dots,a_n\rangle$ be an increasing sequence of points on the real line.
Write a linear time algorithm that determines a unit length interval that contains maximum
number of points of $S$.

\qhead{10}
Let $S$ be a set of size $n$ and $x$ be a number. The sum of the elements of the set $S$ is $m$.
Write an efficient algorithm to determine the subset of $S$ whose sum is closest to $x$ but does
not exceed $x$.

\qhead{10}
A pair $(i,j)$ is called a \emph{significant pair} if $i<j$ and $A[i]>2A[j]$. Design an
$\Oh{n\log n}$-time algorithm to count the number of significant pairs in an array.

\qhead{10}
Let $S$ be a set of $n$ distinct numbers and $k$ be a positive integer such that $k\le n$. Write
an $\Oh n$-time algorithm that determines the $k$ numbers in $S$ that are nearest to the median
of $S$.

%==================================================================
\section{Solutions to the 2022 Paper}

\setcounter{pq}{0}

\qhead{5+5}
\pt{a} \textbf{Counting sort with four counters.}
\begin{enumerate}
  \item Initialise $C[1..4]\leftarrow0$.
  \item One pass: for each $i$, increment $C[A[i]]$.
  \item One pass: write $C[1]$ copies of $1$, then $C[2]$ copies of $2$, then $C[3]$ copies of
        $3$, then $C[4]$ copies of $4$, overwriting $A$ from the front.
\end{enumerate}
\emph{Correctness.} The output is non-decreasing by construction, and it is a permutation of the
input because each value $v$ is written exactly $C[v]=\#\{i:A[i]=v\}$ times.

\emph{Complexity.} Two linear passes plus $\Oh1$ counter work: $\boxed{\Oh n}$ time, $\Oh1$ extra
space (four counters).

\begin{trap}{``but sorting is $\Om{n\log n}$!''}
No contradiction. The $\Om{n\log n}$ bound of \S6 applies only to \textbf{comparison-based}
algorithms. Counting sort never compares two array elements --- it uses their \emph{values as
indices}. Say this sentence explicitly; it is almost certainly worth a mark.
\end{trap}

\pt{b} \textbf{Two-pointer (Lomuto) partition on the $b$-field.}
\begin{enumerate}
  \item $i\leftarrow0$.
  \item For $j=0,1,\dots,n-1$: if $b_j=0$, swap the tuples at positions $i$ and $j$, then
        $i\leftarrow i+1$.
\end{enumerate}
(If the array literally stores the flattened sequence $a_1,b_1,a_2,b_2,\dots$ in $2n$ cells, a
``swap'' is a swap of two adjacent-cell blocks --- still $\Oh1$ per operation.)

\emph{Loop invariant.} Before each iteration, positions $[0,i)$ hold exactly the tuples with
$b=0$ seen so far, and positions $[i,j)$ hold exactly those with $b=1$ seen so far.

\emph{Correctness.} If $b_j=1$ the invariant extends trivially (the region of $1$s grows by one).
If $b_j=0$, the swap sends the tuple to position $i$ --- the first slot of the $1$s block --- and
moves that displaced $1$-tuple to position $j$, which is the new end of the $1$s block; then $i$
advances. Both blocks stay correct. On termination $j=n$, so $[0,i)$ is all the $b=0$ tuples and
$[i,n)$ all the $b=1$ tuples.

\emph{Complexity.} $n$ iterations, each $\Oh1$: $\boxed{\Oh n}$ time, $\Oh1$ space (two indices).

\emph{Remark.} This is exactly the partition primitive of quicksort/quickselect, and the same
idea as the \textsc{Partition} step of median-of-medians (\S3). The output is \textbf{not
stable}; if stability were demanded, $\Oh1$ space and linear time together become genuinely hard.

\qhead{7+3}
\pt{a} \textbf{Andrew's monotone chain.}
\begin{enumerate}
  \item Sort the points lexicographically by $(x,y)$.
  \item \textbf{Lower hull:} scan left to right with a stack; before pushing $p$, while the top
        two stack points $q,r$ and $p$ fail to make a counter-clockwise turn --- i.e.\ the cross
        product $(r-q)\times(p-q)\le0$ --- pop the top.
  \item \textbf{Upper hull:} run the same scan on the points in reverse order.
  \item Concatenate the two chains, dropping the two duplicated endpoints.
\end{enumerate}
(Graham scan is an equally acceptable answer: pick the lowest point, sort the rest by polar
angle about it, then run the same left-turn stack scan.)

\emph{Correctness.} A point is a vertex of the lower hull exactly when it makes a strict left turn
with its two hull neighbours. The stack maintains the invariant ``\emph{the stack currently holds
the lower hull of the points scanned so far}''. Popping removes precisely those points that the
new point $p$ renders non-convex, and a popped point can never return: it lies strictly below the
segment joining two points that remain, so it is inside the hull of the points already seen, and
adding more points only enlarges that hull. Doing the scan in both directions yields the lower and
upper chains, whose concatenation is the full hull.

\emph{Complexity.} Sorting is $\Oh{n\log n}$. In each scan every point is pushed once and popped
at most once, so each scan is $\Oh n$ amortised. Total $\boxed{\Oh{n\log n}}$ --- which by part
(b) is optimal.

\pt{b} \textbf{Lower bound by reduction from sorting.}

Let $\mathcal H$ be any algorithm computing the convex hull (as a cyclically ordered list of
vertices) in time $T(n)$. Given arbitrary reals $x_1,\dots,x_n$ to sort:
\begin{enumerate}
  \item In $\Oh n$ time build the point set $P=\{(x_i,\,x_i^2):i=1,\dots,n\}$, i.e.\ lift every
        number onto the parabola $y=x^2$.
  \item Run $\mathcal H$ on $P$, costing $T(n)$.
  \item In $\Oh n$ time, read the output starting from the leftmost vertex and walking along the
        lower chain, outputting the $x$-coordinates.
\end{enumerate}
\emph{Why step 3 works.} The parabola $y=x^2$ is \textbf{strictly convex}, so \emph{every} point
of $P$ is a vertex of the hull --- no point lies inside or on a chord joining two others. The hull
is therefore a convex polygon on all $n$ points, and traversing its boundary from the leftmost
vertex along the lower chain visits them in strictly increasing order of $x$. So step 3 emits
$x_1,\dots,x_n$ sorted.

Hence sorting $n$ reals is possible in $T(n)+\Oh n$ time. But comparison-based sorting requires
$\Om{n\log n}$ (\S6, decision-tree argument), so
\[
  T(n)+\Oh n=\Om{n\log n}\quad\Longrightarrow\quad \boxed{T(n)=\Om{n\log n}} . \qquad\square
\]

\begin{facts}{the reduction template, in three lines}
To prove problem $\Pi$ needs $\Om{n\log n}$: \textbf{(i)} transform a sorting instance into a
$\Pi$ instance in $\Oh n$; \textbf{(ii)} show a solution to $\Pi$ can be converted back into the
sorted order in $\Oh n$; \textbf{(iii)} conclude $T_\Pi(n)+\Oh n=\Om{n\log n}$. State all three
steps --- marks are awarded per step, and step (ii) is the one people omit.
\end{facts}

\qhead{5+5}
\pt{a} The statement is only meaningful once the array is random, so take $A$ to be a
\textbf{uniformly random permutation} of $n$ distinct numbers.

For each pair $i<j$ define the indicator
\[
  X_{ij}=\mathbf 1\{A[i]>A[j]\},\qquad\text{so}\qquad X=\sum_{1\le i<j\le n}X_{ij}.
\]
Fix such a pair. In a uniformly random permutation the two relative orders of the values at
positions $i$ and $j$ are equally likely, so $\Prob(A[i]>A[j])=\tfrac12$ and
$\E(X_{ij})=\tfrac12$.

Linearity of expectation needs \emph{no} independence (and indeed the $X_{ij}$ are highly
dependent), so
\[
  \E(X)=\sum_{i<j}\E(X_{ij})=\binom n2\cdot\frac12=\boxed{\frac{n(n-1)}{4}} .
\]
\emph{Sanity checks.} $n=2$ gives $\tfrac12$ ($0$ or $1$ inversion, equally likely) \checkmark;
the maximum possible is $\binom n2$, attained by the reversed array, and the mean sits at exactly
half of it, as symmetry demands (reversing a permutation maps $X\mapsto\binom n2-X$).

\pt{b} \textbf{Claim.} If $a>b^k$ then $T(n)=\Th{n^{\log_b a}}$.

Take $n=b^L$, so $L=\log_b n$. Unrolling the recurrence $L$ times:
\[
  T(n)=\sum_{i=0}^{L-1}a^i\,c\!\left(\frac{n}{b^i}\right)^{\!k}+a^L\,T(1)
      =c\,n^k\sum_{i=0}^{L-1}\left(\frac{a}{b^k}\right)^{\!i}+d\,a^L .
\]
Note $a^L=a^{\log_b n}=n^{\log_b a}$. Write $r=a/b^k$; the hypothesis $a>b^k$ gives $r>1$, so the
geometric sum is dominated by its \emph{largest} (last) term:
\[
  \sum_{i=0}^{L-1}r^i=\frac{r^L-1}{r-1}=\Th{r^L},
  \qquad
  r^L=\left(\frac{a}{b^k}\right)^{\log_b n}=\frac{n^{\log_b a}}{n^{k}} .
\]
Substituting,
\[
  c\,n^k\cdot\Th{\frac{n^{\log_b a}}{n^k}}=\Th{n^{\log_b a}},
\]
and the second term $d\,n^{\log_b a}$ is of the same order. Hence
\[
  \boxed{T(n)=\Th{n^{\log_b a}} .}
\]
(For general $n$, monotonicity of $T$ --- given in the question --- lets you sandwich $n$ between
consecutive powers of $b$, and $\lceil\cdot\rceil$ effects change nothing since
$b^{\log_b a}$ is a constant factor.) \qed

\emph{Interpretation.} $r>1$ means the total work \emph{grows} geometrically down the recursion
tree, so the bottom level --- the $a^L=n^{\log_b a}$ leaves --- dominates. That one sentence is
the whole theorem, and is worth writing.

\qhead{10}
\textbf{Algorithm --- two-pointer sweep.}
\begin{enumerate}
  \item $j\leftarrow1$, $\mathit{best}\leftarrow0$, $\mathit{bestStart}\leftarrow1$.
  \item For $i=1,2,\dots,n$:
        \begin{enumerate}
          \item While $j\le n$ and $a_j\le a_i+1$: $j\leftarrow j+1$.
          \item Now the interval $[a_i,\,a_i+1]$ contains exactly $a_i,a_{i+1},\dots,a_{j-1}$, so
                its count is $j-i$. If $j-i>\mathit{best}$, set $\mathit{best}\leftarrow j-i$ and
                $\mathit{bestStart}\leftarrow i$.
        \end{enumerate}
  \item Return the interval $[a_{\mathit{bestStart}},\,a_{\mathit{bestStart}}+1]$ and the count
        $\mathit{best}$.
\end{enumerate}

\begin{correct}
\textbf{Step 1: it suffices to consider intervals whose left endpoint is a point of $S$.}
Let $I=[t,t+1]$ be any unit interval and let $a_i$ be the \emph{smallest} point of $S$ inside
$I$ (if $I$ contains no point, it is not a candidate for the maximum unless the maximum is $0$).
Then $a_i\ge t$, so $a_i+1\ge t+1$, and every point of $S\cap I$ lies in $[a_i,\,t+1]\subseteq
[a_i,\,a_i+1]$. Hence
\[
  |S\cap I|\ \le\ \bigl|S\cap[a_i,a_i+1]\bigr| ,
\]
so no interval beats the best one of the restricted form. Since the algorithm examines
$[a_i,a_i+1]$ for \emph{every} $i$, it examines an optimal interval.

\textbf{Step 2: the count computed is correct.} $S$ is increasing, so the points in
$[a_i,a_i+1]$ are consecutive, starting at $a_i$. The inner loop advances $j$ to the first index
with $a_j>a_i+1$, so the contained points are exactly $a_i,\dots,a_{j-1}$, a count of $j-i$.

\textbf{Step 3: $j$ never needs to move backwards.} $a_i+1$ is increasing in $i$, so the first
index exceeding it is non-decreasing in $i$. Hence carrying $j$ over from the previous iteration
--- rather than restarting it --- is valid. \qed
\end{correct}

\begin{cplx}
$i$ increases $n$ times in total; $j$ increases at most $n$ times \emph{in total across the whole
run} (it never decreases and is capped at $n+1$). Every other operation is $\Oh1$. So the total
is $\boxed{\Oh n}$, with $\Oh1$ extra space.
\end{cplx}

\begin{trap}{do not restart the inner pointer}
The naive ``for each $i$, scan forward from $i$'' is $\Oh{n^2}$ in the worst case (e.g.\ all $n$
points inside one unit interval). The linear bound comes \emph{entirely} from the observation in
Step 3 that $j$ is monotone. Say it explicitly --- that is where the marks are.
\end{trap}

\qhead{10}
\textbf{Assumption.} The elements of $S$ are positive (otherwise shift them; the problem's mention
of the total $m$ signals this reading). Write $S=\{s_1,\dots,s_n\}$ and set $B=\min(x,m)$ --- if
$x\ge m$ the answer is all of $S$, so we may assume the target is capped at $m$.

\textbf{Algorithm --- subset-sum dynamic programming.}
\begin{enumerate}
  \item Boolean table $R[0..B]$, with $R[0]=\text{true}$ and $R[s]=\text{false}$ otherwise. (Also
        keep $\mathit{take}[s]$, the index of the element used to first reach $s$, for recovery.)
  \item For each element $s_t\in S$, for $s=B$ \textbf{down to} $s_t$:
        if $R[s-s_t]$ and not $R[s]$, set $R[s]\leftarrow\text{true}$ and
        $\mathit{take}[s]\leftarrow t$.
  \item Let $s^\star=\max\{s\le B:R[s]=\text{true}\}$ --- scan $R$ downwards from $B$.
  \item Recover the subset by repeatedly following $\mathit{take}$: from $s^\star$, output element
        $\mathit{take}[s^\star]$ and move to $s^\star-s_{\mathit{take}[s^\star]}$, until reaching
        $0$.
\end{enumerate}

\begin{correct}
\textbf{Invariant.} After processing the first $t$ elements, $R[s]$ is true exactly when some
subset of $\{s_1,\dots,s_t\}$ sums to $s$.

\emph{Base:} with no elements only $s=0$ is reachable. \emph{Step:} a subset of the first $t$
elements summing to $s$ either omits $s_t$ (so $s$ was already reachable) or includes it (so
$s-s_t$ was reachable using the first $t-1$), which is precisely the update rule.

\textbf{Why the loop runs downwards.} Iterating $s$ from high to low guarantees that
$R[s-s_t]$ still refers to the table \emph{before} $s_t$ was processed. Ascending order would let
$s_t$ be used repeatedly, solving the \emph{unbounded} knapsack instead.

Since every achievable sum $\le B$ is marked, $s^\star$ is the largest achievable sum not
exceeding $x$ --- i.e.\ the closest from below. \qed
\end{correct}

\begin{cplx}
$n$ elements, each sweeping a table of size $B+1$:
\[
  \boxed{\Oh{n\,B}=\Oh{n\min(x,m)}\ \le\ \Oh{nm}}
\]
time, $\Oh{B}$ space (the $\mathit{take}$ array costs the same). Recovery is $\Oh n$.
\end{cplx}

\begin{facts}{say the word ``pseudo-polynomial''}
Subset-sum is \textbf{NP-hard}, so no algorithm polynomial in the \emph{input length} is expected;
$\Oh{nm}$ is polynomial in the \emph{value} $m$, not in its bit-length $\Oh{\log m}$. That is what
``pseudo-polynomial'' means, and it is exactly why the question says ``\emph{efficient}'' rather
than naming a complexity. Contrast with the brute force over all $2^n$ subsets.
\end{facts}

\qhead{10}
\textbf{Algorithm --- modified merge sort.} \code{Count}$(A[\ell..r])$ returns the number of
significant pairs inside the range \emph{and} leaves the range sorted ascending.
\begin{enumerate}
  \item If $\ell\ge r$, return $0$.
  \item $m\leftarrow\floor{(\ell+r)/2}$; recursively
        $c_L\leftarrow\code{Count}(A[\ell..m])$ and $c_R\leftarrow\code{Count}(A[m{+}1..r])$.
        Both halves are now sorted.
  \item \textbf{Cross-count.} Let $L=A[\ell..m]$ and $R=A[m{+}1..r]$. Walk a pointer $q$ over $R$
        in increasing order and maintain a pointer $p$ into $L$ at the first position with
        $L[p]>2R[q]$. Add $|L|-p$ to a counter $c_X$. Because $2R[q]$ increases with $q$, the
        pointer $p$ only moves \textbf{forward}.
  \item Merge $L$ and $R$ into $A[\ell..r]$ in the standard way.
  \item Return $c_L+c_R+c_X$.
\end{enumerate}

\begin{correct}
\textbf{Every significant pair is counted exactly once.} Merge sort splits by \textbf{index}, so
for a pair $(i,j)$ with $i<j$ exactly one of three things holds: both indices lie in the left
half (counted by the recursive call on $L$), both in the right half (counted by the call on $R$),
or $i$ is in the left half and $j$ in the right half (counted in step 3). The third case is
correct because \emph{every} left-half index is smaller than \emph{every} right-half index, so the
condition $i<j$ is automatic and only $A[i]>2A[j]$ needs testing.

\textbf{The cross-count is correct.} With $L$ sorted ascending, $\{p:L[p]>2R[q]\}$ is an upward-
closed set, so its size is $|L|-p$ where $p$ is the first such position --- which is what the
pointer tracks.

\textbf{Sorting the halves is harmless.} Whether a pair is significant depends only on the two
\emph{values} and on which half each index came from, never on the order within a half. So
reordering inside a half after its own pairs have been counted changes nothing.

\textbf{Monotonicity is what makes the sweep legal.} $R$ is scanned in increasing order and
$z\mapsto2z$ is increasing, so the threshold $2R[q]$ increases and $p$ never needs to retreat.
\qed
\end{correct}

\begin{cplx}
Steps 3 and 4 are each one linear pass, so
\[
  T(n)=2T(n/2)+\Oh n \quad\Longrightarrow\quad \boxed{\Oh{n\log n}}
\]
by the Master Theorem, Case 2 ($a=b=2$, $n^{\log_b a}=n$, $f(n)=\Th n$).
\end{cplx}

\begin{facts}{one skeleton, many questions}
Identical structure solves: counting inversions ($A[i]>A[j]$), this problem ($A[i]>2A[j]$), and
practice Q5 ($A[i]>3A[j]+5$). \textbf{The only requirement is that the threshold function
$g(A[j])$ be monotone increasing} --- that is what permits the single forward pointer. If it were
not monotone you would need a balanced BST or BIT per merge, giving $\Oh{n\log^2n}$.
\end{facts}

\qhead{10}
\textbf{Algorithm --- two applications of median-of-medians.}
\begin{enumerate}
  \item Find the median $M$ of $S$ by median-of-medians (\S3). \hfill $\Oh n$
  \item Form the array of distances $d_i=|s_i-M|$ for every $s_i\in S$. \hfill $\Oh n$
  \item Find the $k$-th smallest value $\delta$ among $d_1,\dots,d_n$, again by
        median-of-medians. \hfill $\Oh n$
  \item Output every $s_i$ with $d_i<\delta$; then top the answer up to exactly $k$ elements using
        elements with $d_i=\delta$. \hfill $\Oh n$
\end{enumerate}

\begin{correct}
Step 3 gives the value $\delta$ such that exactly $k$ of the $d_i$ are $\le\delta$ when counted
with multiplicity. Step 4 outputs a set $K$ of size exactly $k$ consisting of elements with the
$k$ smallest distances: every element with $d_i<\delta$ must be in any optimal answer, and the
remaining slots are filled with elements at distance exactly $\delta$, which are interchangeable.

Hence for any $s\in K$ and $s'\notin K$ we have $|s-M|\le|s'-M|$, which is precisely the
statement that $K$ consists of $k$ numbers nearest the median. (With $S$ a set of \emph{distinct}
numbers, ties $d_i=d_j$ can still occur --- one element on each side of $M$ --- so step 4's
tie-handling is genuinely needed, and any choice is optimal.) \qed
\end{correct}

\begin{cplx}
Four phases, each $\Oh n$ in the worst case:
\[
  \boxed{\Oh n} .
\]
The whole point is that \emph{both} selections use the worst-case-linear algorithm. Sorting $S$
($\Oh{n\log n}$) or sorting the $d_i$ would miss the bound.
\end{cplx}

\begin{trap}{do not sort, and do not use randomised quickselect}
Two tempting wrong answers:
\begin{itemize}
  \item ``Sort $S$, take the $k$ values around the middle'': $\Oh{n\log n}$, fails the required
        $\Oh n$.
  \item ``Use randomised quickselect'': that is $\Oh n$ \emph{expected}, but $\Oh{n^2}$ in the
        worst case. The question says $\Oh n$-time, so you need the deterministic
        median-of-medians guarantee. Name it.
\end{itemize}
Also note the answer is \emph{not} simply ``the $k$ elements around position $n/2$ in sorted
order'' unless you have already sorted --- and you must not.
\end{trap}
